For fixed \(\chi\), every observational cell of \(C_d=(A,X,Z_{M_d})\), for \(d\in\{d_{\mathrm p},d_{\mathrm e}\}\), has probability bounded below by a constant \(c_\chi>0\), uniformly over \(0<\epsilon\le10^{-3}\).  Indeed, \(X\) is constant, the observational treatment kernel lies in \([3/8,5/8]\), every binary proxy kernel used here has both masses bounded away from zero, and the finite mixture over \(U\) preserves a positive lower bound.  Also \(|Y|=1\).  These facts use only the finite positive cells in Definition~\(\mathrm{def{:}binary\text{-}witness}\).

For any bounded cell indicator or cell-marked outcome \(W_i\), independence gives
\[
 \mathbb E\left[\left\{m^{-1}\sum_{i=1}^m(W_i-\mathbb EW_i)\right\}^2\right]
 =m^{-1}\operatorname{Var}(W_1)\le m^{-1}.
\]
Consequently, for every \(M>0\), the probability that the absolute average exceeds \(Mm^{-1/2}\) is at most \(M^{-2}\).  Using the finite number of cells shows that every empirical cell mass differs from its population mass by \(O_P(m^{-1/2})\), uniformly over the two designs and folds.  Hence all observational training-cell counts are at least \(c_\chi m/4\) with probability tending to one.  The ratio identity for a cell mean then gives, uniformly over the two designs, two folds, and finitely many cells,
\[
 \widehat m_{d,m}^{(-\ell)}(c)-m_d(c)=O_P(m^{-1/2}),
 \qquad
 m_d(c)=\mathbb E_P(Y\mid C_d=c,G=O).                                         \tag{5}
\]
The zero-denominator convention is therefore asymptotically inactive.

Conditional on a training half, the corresponding evaluation observations are independent of its fitted cell means, and every squared residual is bounded by four.  The same conditional second-moment calculation gives
\[
 \widehat R_{d,m}^A-
 \mathbb E_P\!\left[
 \{Y-\widehat m_{d,m}^{(-\ell)}(C_d)\}^2\mid G=O,
 \text{training half}\right]
 =O_P(m^{-1/2})                                                               \tag{6}
\]
after averaging the two folds; the random observational denominator differs from its positive expectation by \(O_P(m^{-1/2})\) and is absorbed in the same order.  The conditional-expectation Pythagorean identity yields
\[
 \mathbb E_P\{(Y-\widehat m(C_d))^2\mid G=O,\widehat m\}
 =R_d^A(P)+\sum_cP(C_d=c\mid G=O)\{\widehat m(c)-m_d(c)\}^2.                 \tag{7}
\]
By (5), the last term is \(O_P(m^{-1})\), uniformly over the finite catalogue.  Equations (6)--(7) prove (1).

The center risk calculation in Lemma~\(\mathrm{lem{:}sharp\text{-}uniform\text{-}binary\text{-}two\text{-}risk\text{-}elbow}\) gives, uniformly in \(m\),
\[
 R_{d_{\mathrm e}}^A(P_{\chi,\epsilon_m})
 -R_{d_{\mathrm p}}^A(P_{\chi,\epsilon_m})
 >\frac{585}{1024}-\frac3{10}
 =\frac{1389}{5120}>0.                                                        \tag{8}
\]
When the maximum error in (1) is less than half this fixed gap, the deterministic tie-broken empirical minimizer is \(d_{\mathrm p}\).  This proves (2).

On the event in (2), the random ratio in (3) equals \(\rho_{\mathrm{pred}}^A(P_{\chi,\epsilon_m})\), while the ratio in (4) equals \(\rho_{\mathrm{pred}}^{A,\mathrm{IKM}}(P_{\chi,\epsilon_m})\).  The same lemma gives the deterministic lower bound
\[
 \rho_{\mathrm{pred}}^A(P_{\chi,\epsilon_m})
 =\rho_{\mathrm{pred}}^{A,\mathrm{IKM}}(P_{\chi,\epsilon_m})
 >\frac{\epsilon_m^{-2}+14}{726}\longrightarrow\infty.                     \tag{9}
\]
Combining (2) and (9) proves convergence in probability in (3)--(4).